\documentclass[10pt,a4paper]{amsart}
\usepackage[margin=19mm]{geometry}
\usepackage{amsmath,amssymb,mathtools}
\usepackage[scr=rsfs,cal=euler]{mathalfa}
\usepackage{xfrac}
\usepackage[unicode,hidelinks,bookmarks=false]{hyperref}
\newcommand{\ie}{i.e.\ }
\newcommand{\cf}{cf.\ }
\newcommand{\resp}{respectively\ }
\newcommand{\Subsection}{\subsection}
\providecommand{\nobreakdash}{\nobreak\hbox{-}}
\setlength{\emergencystretch}{2em}
\multlinegap0pt
\usepackage{bm}
\usepackage{bbm}
\usepackage{dsfont}
\usepackage{xspace}
\usepackage{cancel}
\usepackage{mathtools}
\usepackage{amsthm}
\makeatletter
\@ifundefined{theorem}{\newtheorem{theorem}{Theorem}}{}
\@ifundefined{lemma}{\newtheorem{lemma}[theorem]{Lemma}}{}
\@ifundefined{proposition}{\newtheorem{proposition}[theorem]{Proposition}}{}
\makeatother
\title[Nonexistence of the metric with positive intermediate curvat]{Nonexistence of the metric with positive intermediate curvatures on manifolds with boundary}
\author{Jingche Chen, Han Hong}
\date{Source: arXiv:2510.13099v1 (CC BY 4.0)}
\begin{document}
\maketitle

\noindent\emph{Source and licence.} Nonexistence of the metric with positive intermediate curvatures on manifolds with boundary. Version arXiv:2510.13099v1 (CC BY 4.0), distributed under the Creative Commons Attribution 4.0 International licence, as recorded in the arXiv licence metadata for this exact version and shown on the abstract page https://arxiv.org/abs/2510.13099v1. Full rights notice: licenses/arxiv-2510.13099-CC-BY-4.0.txt.\par\smallskip

\noindent\textit{Editorial scope.} Counted unit: the Lemma whose source label is \texttt{foliation} in the article (source0.tex lines 813--816, source label \texttt{foliation}), delivered together with the article's own complete original proof of it (source0.tex lines 817--851, one \texttt{proof} environment of 35 lines). No mathematics is written, repaired or re-derived: statement and proof are the article's own text line for line. Required context retained uncounted: main theorem 1 Introduction (lines 528--535); rigidity (lines 787--807). Every definition, display and label of those lines is retained unchanged. Omitted from this package and not redistributed: 1--527, 536--786, 808--812, 852--1274. Nothing inside a retained excerpt is omitted or altered: every retained body is delivered exactly as the article's lines are written, so no omission receipt of any kind accompanies this package. Citations: the retained excerpts carry no citation command at all, so no bibliography entry of the article is transcribed and no reference is redistributed. Reference adaptation: none was needed and none was made; every reference inside the delivered unit resolves inside it, so no unresolved reference or citation is produced. Printed slips kept literal: no printed word, symbol, hypothesis or number of the counted statement or of its proof was altered. Layout and class: the article's own class is \texttt{amsart} with packages including amsmath, amssymb, amsthm, hyperref, cleveref, mathrsfs, mathtools, slashed, tikz-cd, enumitem; the wrapper is the lane's pinned \texttt{amsart} editorial wrapper at 10pt on A4, which restates the theorem-like environments the retained text needs and adds the article's own macro definitions that the retained text needs (none), with extra wrapper packages (bm, bbm, dsfont, xspace, cancel, mathtools). The delivered numbering is the unit's own. Assets: no retained span contains a figure environment, a graphics-inclusion command, a PDF-page inclusion, a table or a TikZ picture; no image file is copied into this package, the package declares no asset dependency and no image file is redistributed with it. Licence: version arXiv:2510.13099v1 of this article is distributed under the Creative Commons Attribution 4.0 International licence, as recorded in the arXiv licence metadata for this exact version and shown on the abstract page https://arxiv.org/abs/2510.13099v1. A full rights notice is delivered at licenses/arxiv-2510.13099-CC-BY-4.0.txt. Characters: every retained excerpt is pure ASCII, so the delivered engine, XeLaTeX with its default Latin Modern text font, reports no missing character.
\begin{theorem}\label{main theorem 1 Introduction}
    Let $(M,\partial M)$ be a compact orientable $n$-dimensional manifold. Suppose that there exist nonzero $  \alpha^{1},...,\alpha^{m}\in H^{1}(M,\mathbb{Z})$ such that %$\alpha^{1}\smile\cdots\smile \alpha^{m}\neq 0$ and
    $(\alpha^{1}\smile\cdots\smile \alpha^{m})|_{\partial M}\neq 0$. Then

    (1) If $2\leq n\leq 7$ and $1\leq m\leq n-1$, there is no Riemmanian metric $g$ on $M$ satisfying $C_{m}> 0$ and $H_{m}\geq 0$, or $C_{m}\geq 0$ and $H_{m}> 0$.  

    (2) If $2\leq n\leq 6$ and $1\leq m\leq n-1$, or $n=7$ and $m\in \{1,5,6\}$, and there exists a metric $g$ on $M$ satisfying $C_{m}\geq 0$ and $H_{m}\geq 0$, then $(M,\partial M)$ is isometrically covered by $(X^{n-m}\times \mathbb{R}^{m},\partial X^{n-m}\times \mathbb{R}^{m})$, where $X$ is a compact manifold with boundary satisfying $\text{Ric}_{X}\geq 0$ and $A^{\partial X}\geq 0$.
\end{theorem}

\begin{proposition}\label{rigidity}
    Under the assumption of (2) of Theorem \ref{main theorem 1 Introduction}, the following is true:

    $(1)$ $\lambda_{k}=0$ for all $ 1\leq k\leq m-1$;

    $(2)$ $\nabla_{\Sigma_{m-1}}\log\rho_{m-1}=0$ along $\Sigma_{m}$;

    $(3)$ $\text{Ric}_{\Sigma_{m-1}}(\nu_{m},\nu_{m})=(\nabla^{2}_{\Sigma_{m-1}}\log \rho_{m-1})(\nu_{m},\nu_{m})$ on $\Sigma_{m}$;

    $(4)$ $C_{m}(\nu_{1},...,\nu_{m})=0$;

    $(5)$ $R_{\Sigma_{m}}=R_{M}|_{\Sigma_{m}}$;

    $(6)$ $A^{\Sigma_{k}}=0$ for all $\ 1\leq k\leq m$ on $\Sigma_{m}$;

    $(7)$ $\langle \nu_{k},\eta\rangle =0$ for all $\ 1\leq k\leq m$ along $\partial \Sigma_m$;
    
    $(8)$ $A^{\partial \Sigma_{m-1}}(\nu_{m},\nu_{m})=0$;

    $(9)$ $H_{m}(\nu_{1},...,\nu_{m})=0$ in $\partial\Sigma_{m-1}\subset \Sigma_{m-1}$.
\end{proposition}

\begin{lemma}\label{foliation}
    If $\Sigma_{m}\subset \Sigma_{m-1}$ satisfies assumptions in Proposition \ref{rigidity}, then we can find a local foliation $\{\Sigma_{m,t}\}_{-\epsilon<t<\epsilon}$ of $\Sigma_{m}$ in $\Sigma_{m-1}$ such that $\Sigma_{m,t}$ is given by the graph over $\Sigma_m$ with graph function $u_t$ along the unit normal $\nu_m$. In other words, we can define a map $f_{t}:\Sigma_{m}\rightarrow \Sigma_{m-1}$ by $f_{t}(x)=\exp_{x}(u_{t}\nu_{m})$ and $\Sigma_{m,t} = f_{t}(\Sigma_{m})$. Moreover, it satisfies $$\Sigma_{m,0}=\Sigma_{m}, \left.\frac{\partial}{\partial t}\right|_{t=0}u_{t}=1, \frac{\partial u_{t}}{\partial t}>0, \int_{\Sigma_{m}}(u_{t}-t)=0, \langle \nu_{m,t},\eta_{t}\rangle=0$$ and $H_{\Sigma_{m,t}}+\langle \nabla_{\Sigma_{m-1}}\log \rho_{m-1},\nu_{m,t}\rangle = \text{const}$ on $\Sigma_{m,t}$ where $\nu_{m,t}$ is the unit normal vector field of $\Sigma_{m,t}$ and $\eta
    _{t}$ is the unit normal vector field of $\partial \Sigma_{m-1}\hookrightarrow \Sigma_{m-1}$ restricted to $\partial \Sigma_{m,t}$.
\end{lemma}

\begin{proof}
    Consider the Banach space $\mathcal{W}=\{(f,g)\in C^{\alpha}(\Sigma_{m})\times C^{\alpha}(\partial \Sigma_{m})|\int_{\Sigma_{m}}f-\int_{\partial \Sigma_{m}}g=0\}$. Given a function $u\in C^{2,\alpha}(\Sigma_{m})$, let $H_{\Sigma_u}$ denote the mean curvature of the graph $\Sigma_u$. Denote $\tilde{H}_{\Sigma_{u}}=H_{\Sigma_{u}}+\langle \nabla_{\Sigma_{m-1}}\log\rho_{m-1},\nu_{u}\rangle$ where $\nu_{u}$ is the unit normal vector field of $\Sigma_{u}$ and $\eta
    _{u}$ is the unit normal vector field of $\partial \Sigma_{m-1}\hookrightarrow \Sigma_{m-1}$ restricted to $\partial \Sigma_{u}$. Notice that $\Sigma_{m}$ is free boundary, then $\eta:=\eta_{u}|_{u=0} = \eta_{\Sigma_{m}}$, where $\eta_{\Sigma_{m}}$ is the outward conormal vector field of $\Sigma_{m-1}$. We can define a map $\Psi:C^{2,\alpha}(\Sigma_{m})\rightarrow\mathcal{W}\times \mathbb{R}$ as the following
    \[
    \Psi(u)=\Big(\big(\tilde{H}_{\Sigma_{u}}-\frac{1}{|\Sigma_{m}|}\int_{\Sigma_{m}}\tilde{H}_{\Sigma_{u}}+\frac{1}{|\Sigma_{m}|}\int_{\partial\Sigma_{m}}\langle \nu_{u},\eta_{u}\rangle,\langle\nu_{u},\eta_{u}\rangle\big),\frac{1}{|\Sigma_{m}|}\int_{\Sigma_{m}}u\Big).
    \]
    It is easy to check that $\Psi(0)=((0,0),0)$. Now we want to compute $D\Psi|_{u=0}$. We notice that
    \[
    \begin{aligned}
    &D\tilde{H}_{\Sigma_{u}}|_{u=0}(w)\\
    =&-\Delta_{\Sigma_{m}}w-|A^{\Sigma_{m}}|^2w-\text{Ric}_{\Sigma_{m-1}}(\nu_{m},\nu_{m})w+(\nabla^{2}_{\Sigma_{m-1}}\log \rho_{m-1})(\nu_{m},\nu_{m})w\\
    =&-\Delta_{\Sigma_{m}}w
    \end{aligned}
    \]
    where we applied $(3)$ and $(6)$ in Proposition \ref{rigidity}. 
    We also compute 
    \[
    \begin{aligned}
    D\langle\nu_{u},\eta_{u}\rangle|_{u=0}(w) 
    &=\langle -\nabla_{\Sigma_{m-1}}w,\eta\rangle + w\langle \nu_{m},\nabla_{\nu_{m}}\eta\rangle\\
    &=-\frac{\partial w}{\partial \eta}+ wA^{\partial \Sigma_{m-1}}(\nu_{m},\nu_{m})\\
    &=-\frac{\partial w}{\partial \eta}
    \end{aligned}
    \]
    where the final equality is due to $A^{\partial \Sigma_{m-1}}(\nu_{m},\nu_{m})=0$ from Proposition \ref{rigidity}. Thus, we have
    \[
    D\Psi|_{u=0}(w) = \Big(\big(-\Delta_{\Sigma_{m}}w,-\frac{\partial w}{\partial\eta}\big),\frac{1}{|\Sigma_{m}|}\int_{\Sigma_{m}}w\Big)
    \]
    where we used the fact that $\int_{\Sigma_{m}}\Delta_{\Sigma_{m}}w=\int_{\partial \Sigma_{m}}\frac{\partial w}{\partial \eta}$ and $\eta = \eta_{\Sigma_{m}}$.
    By the uniqueness and existence of solutions to the elliptic equation under Nuemann boundary condition, we obtain that $D\Psi|_{u=0}$ is injective and surjective. According to the inverse function theorem, there exists a small $\epsilon>0$ such that for any $t\in (-\epsilon,\epsilon)$, there exists $u_{t}$ satisfying $\Psi(u_{t})=((0,0),t)$. Differentiating the equation $\Psi(u_{t})=((0,0),t)$, we have
    \[
    \Big(\big(-\Delta_{\Sigma_{m}}(\left.\frac{\partial }{\partial t}\right|_{t=0}u_{t}),-\frac{\partial}{\partial\eta}(\left.\frac{\partial}{\partial t}\right|_{t=0}u_{t})\big),\frac{1}{|\Sigma_{m}|}\int_{\Sigma_{m}}(\left.\frac{\partial}{\partial t}\right|_{t=0}u_{t})\Big)=(0,0,1).
    \]
    Then $\left.\frac{\partial}{\partial t}\right|_{t=0}u_{t}=1$. Taking $\epsilon$ smaller, we get $\partial_{t}u_{t}>0$. 
\end{proof}

\end{document}
